\documentclass[10pt,a4paper]{article}
\usepackage[margin=19mm]{geometry}
\usepackage[T1]{fontenc}
\usepackage{lmodern}
\usepackage{microtype}
\usepackage{amsmath,amssymb,amsthm,mathtools}
\usepackage{graphicx}
\usepackage{xcolor}
\usepackage{enumitem}
\usepackage[hidelinks]{hyperref}

\definecolor{navy}{RGB}{21,55,94}
\definecolor{softblue}{RGB}{238,245,252}
\newtheorem{theorem}{Theorem}
\newtheorem{corollary}{Corollary}
\newtheorem{lemma}{Lemma}
\newtheorem{remark}{Remark}
\newcommand{\C}{\mathbb C}
\newcommand{\Mat}{\operatorname{Mat}}
\newcommand{\rk}{\operatorname{rank}}
\newcommand{\spec}{\operatorname{spec}}

\title{\vspace{-1.4em}\textbf{Rank stability for normal matrices of\\bounded spectral complexity}}
\author{Candidate partial-result packet for arXiv:1708.05338}
\date{17 August 2026}

\begin{document}
\maketitle
\vspace{-1.4em}

\begin{center}
\fcolorbox{navy}{softblue}{\begin{minipage}{0.91\linewidth}
\textbf{Outcome.} Fix the tuple length $n$ and a spectral-complexity bound
$M$.  Pairwise almost commuting normal matrices, each having at most $M$
distinct eigenvalues, are uniformly close in normalized rank to an exactly
commuting tuple.  Invertibility is unnecessary.  The conclusion persists if
all but a sufficiently small rank fraction of each spectrum is supported on
at most $M$ values.
\end{minipage}}
\end{center}

\section*{1. The source question and the partial answer}

For $X,Y\in\Mat_d(\C)$ write
\[
 d_{\rm rk}(X,Y)=\frac{\rk(X-Y)}d,
 \qquad \|[X,Y]\|_{\rm rk}=\frac{\rk(XY-YX)}d.
\]
Elek--Grabowski prove stability when every matrix is self-adjoint or unitary,
and ask for arbitrary normal invertible matrices.

\begin{figure}[h]
\centering
\includegraphics[width=0.96\linewidth]{figures/open_problem_crop.png}
\caption{The exact questions in Remark 2, source PDF page 3.}
\end{figure}

\begin{theorem}[Bounded spectral complexity]\label{thm:main}
For every $n,M\ge1$ and $\varepsilon>0$ there is $\delta>0$, independent of
$d$, with the following property.  If $A_1,\ldots,A_n\in\Mat_d(\C)$ are
normal, $|\spec(A_i)|\le M$, and
\[
 \|[A_i,A_j]\|_{\rm rk}\le\delta \qquad(1\le i,j\le n),
\]
then there are pairwise commuting $B_1,\ldots,B_n\in\Mat_d(\C)$ such that
$d_{\rm rk}(A_i,B_i)\le\varepsilon$ for every $i$.
\end{theorem}

\paragraph{Proof intuition.}
On a finite spectrum, complex conjugation is exactly a polynomial: if $A$ is
normal with at most $M$ spectral values, then $A^*=p(A)$ for some
$\deg p<M$.  A telescoping identity bounds the rank cost of passing from $A$
to $p(A)$ solely in terms of $M$.  Thus adjoining all adjoints produces the
star-closed almost commuting tuple required by the source's general theorem.

\section*{2. Polynomial commutators}

\begin{lemma}\label{lem:poly}
Let $p\in\C[z]$ have degree at most $m$.  For square matrices $X,Y$,
\[
 \rk[p(X),Y]\le \frac{m(m+1)}2\,\rk[X,Y].
\]
\end{lemma}

\begin{proof}
For $\ell\ge1$,
\[
 [X^\ell,Y]=\sum_{q=0}^{\ell-1}X^q[X,Y]X^{\ell-1-q},
\]
so its rank is at most $\ell\rk[X,Y]$.  Expand $p$ and use subadditivity of
rank.  Scalar coefficients do not affect the estimate.
\end{proof}

\begin{proof}[Proof of Theorem~\ref{thm:main}]
Put $K=\max\{1,M(M-1)/2\}$.  Apply Theorem 6 of Elek--Grabowski to tuples
of length $2n$ and tolerance $\varepsilon$, and call its commutator tolerance
$\eta$; choose $\delta=\eta/K^2$.

By the spectral theorem, for each $i$ there is a polynomial $p_i$ of degree
at most $M-1$ such that
\[
 p_i(\lambda)=\overline\lambda\quad(\lambda\in\spec(A_i)),
 \qquad\text{hence}\qquad A_i^*=p_i(A_i).
\]
Lemma~\ref{lem:poly} gives
\begin{align*}
 \rk[A_i^*,A_j]&\le K\rk[A_i,A_j],\\
 \rk[A_i^*,A_j^*]&\le K^2\rk[A_i,A_j].
\end{align*}
The second inequality is obtained by applying the lemma once to each
polynomial.  Thus
$(A_1,\ldots,A_n,A_1^*,\ldots,A_n^*)$ is a star-closed
$\eta$-commuting tuple in normalized rank.  The source's Theorem 6 supplies
a commuting $2n$-tuple within $\varepsilon$; its first $n$ entries are the
required $B_i$.  No step uses invertibility.
\end{proof}

\section*{3. A robust spectral-concentration upgrade}

\begin{corollary}\label{cor:concentration}
For every $n,M\ge1$ and $\varepsilon>0$ there exist $\delta,\kappa>0$,
independent of $d$, such that the following holds.  Suppose the normal
matrices $A_1,\ldots,A_n\in\Mat_d(\C)$ satisfy
$\|[A_i,A_j]\|_{\rm rk}\le\delta$, and for each $i$ some set
$S_i\subset\spec(A_i)$ of at most $M$ values has total eigenspace dimension
at least $(1-\kappa)d$.  Then the conclusion of Theorem~\ref{thm:main} holds.
\end{corollary}

\begin{proof}
Use Theorem~\ref{thm:main} with tolerance $\varepsilon/2$, and denote its
input tolerance by $\eta$.  Choose
\[
 \delta\le\eta/2,
 \qquad
 \kappa\le\min\{1/2,\varepsilon/2,\eta/8\}.
\]
In an orthonormal eigenbasis for $A_i$, replace every eigenvalue outside
$S_i$ by one fixed member of $S_i$; call the resulting normal matrix $C_i$.
Then $|\spec(C_i)|\le M$ and
$d_{\rm rk}(A_i,C_i)\le\kappa$.  If $D_i=C_i-A_i$, then
\[
 [C_i,C_j]-[A_i,A_j]=[D_i,C_j]+[A_i,D_j],
\]
and consequently
\[
 \|[C_i,C_j]\|_{\rm rk}
 \le\delta+2d_{\rm rk}(A_i,C_i)+2d_{\rm rk}(A_j,C_j)
 \le\eta.
\]
Apply Theorem~\ref{thm:main} to the $C_i$ and use the rank triangle
inequality.
\end{proof}

\section*{4. Why this does not yet prove the full conjecture}

The dependence on $M$ is genuine in this argument: interpolating the adjoint
on an unrestricted finite spectrum can require degree $d-1$, destroying the
dimension-free commutator modulus.  Nor can the source's one-line
unitary/self-adjoint reduction simply be copied.  Indeed, let
\[
 v=(1,2,3,4)^T,\quad U=I-\frac{2vv^*}{v^*v},
\]
\[
 A=\operatorname{diag}(0,1,-5i,2+i),\qquad
 B=U\operatorname{diag}(0,1+i,2-i,3+2i)U^*.
\]
Both $A$ and $B$ are normal, but exact row reduction gives
\[
 \rk[A,B]=2,\qquad \rk[A^*,B]=4.
\]
Thus even equality of the two commutator ranks fails for normal pairs.  This
example obstructs the simplest rank-Fuglede shortcut; it is not a
counterexample to stability.

The theorem also does not answer the separate same-type question: Theorem 6
produces arbitrary commuting approximants, with no guarantee that they remain
normal or invertible.

\paragraph{Upgrade audit.}
Eight focused stages were completed: direct star-closure; later-literature
audit; exact rank-Fuglede stress test; spectral interpolation; quantitative
Horner control; star-closed completion; spectral-concentration upgrade; and
unrestricted/same-type obstruction analysis.  The successful route closes
every uniformly bounded spectral-complexity class.  The remaining barrier is
dimension-free control of adjoints for spectra whose cardinality grows with
$d$.

\paragraph{Novelty boundary.}
A 2024 primary source still records the unrestricted spectral-restriction-free
problem as open.  Bounded searches through 17 August 2026 found no paper
stating this partial theorem.  Expert review and a broader priority search are
required before dissemination.

\begin{thebibliography}{9}\small
\bibitem{EG}
G.~Elek and L.~Grabowski,
\emph{Almost commuting matrices with respect to the rank metric},
arXiv:1708.05338 (2017; v3, 2021).

\bibitem{BBG}
T.~Bauer, G.~Blachar, and B.~Greenfeld,
\emph{Rank-stability of polynomial equations},
arXiv:2401.04676 (2024).
\end{thebibliography}

\end{document}
